% Related work — covers existing DSP Eurorack platforms, interpolation literature, and reverb lineage
% Placement: end of Section 2 (Background and Fundamentals), before System Design

\subsection{Related Work}
\label{sec:related-work}

\subsubsection{Open DSP Eurorack Platforms}

Several open-source and semi-open hardware platforms exist that pursue a similar goal:
bringing programmable DSP capability to the Eurorack format.

\paragraph{Daisy Seed \cite{ElectroSmithHardware}} (Electro-Smith, 2020) is perhaps the most widely adopted platform in this space.
It is based on an STM32H750 microcontroller and is sold as a ready-made module core,
with a C++ DSP library (libDaisy) and a visual patching environment (Max/MSP via RNBO) targeting it.
% TODO: cite Electro-Smith Daisy documentation
Its main advantage is its low barrier to entry; a developer can target the platform without
designing any hardware.
The trade-off is a fixed form factor and pin mapping, which constrains the possible interface design.
This thesis takes the opposite approach: the hardware is designed from scratch around the
application's specific requirements, at the cost of higher initial development effort.

\paragraph{Rebel Technology OWL \cite{Genius2022}} is a further example of a programmable Eurorack DSP module,
targeting a similar embedded ARM platform.

\paragraph{Mutable Instruments Eurorack Modules}
The open-source Eurorack modules developed by Mutable Instruments \cite{pichenettesPichenettesEurorack2026} served as an important reference throughout this work.
In particular, the Clouds texture synthesizer was studied as a design reference for several aspects of this project.
The CV input stage topology — specifically the use of a \qty{-10}{\volt} bias reference in an inverting summing amplifier to shift bipolar CV signals into a unipolar ADC input range — is directly inspired by the Clouds hardware design.
The calibration scheme and the phase accumulator approach used in the oscillator and playback engine are likewise informed by the Clouds firmware implementation.
The Mutable Instruments codebase is published under open-source licences and the hardware under a Creative Commons licence, making it a legitimate and well-documented reference for derivative design work.

None of the above platforms are specifically designed for tape-style sample playback, which is the primary application explored in this thesis.
Furthermore, all of the above are module products rather than custom hardware designs —
this work contributes a documented, open hardware platform that can be adapted and extended.
